%% ECON 282E -- Session 5, Deck A4: Linear Rational-Expectations Systems
%% Source: Quant_Macro Lec.9 frames 25-30 (Blanchard-Kahn), PLUS material absent
%% from that deck entirely -- log-linearization, undetermined coefficients, the
%% matrix quadratic, QZ/Klein, Uhlig, Sims, Dynare.  Those come from
%% Lecture_SM_5_Appendix_Linearization.pdf pp.6-19 and
%% Lecture_SM_5_Perturbation_I.pdf pp.20, 30 (Fernandez-Villaverde &
%% Guerron-Quintana), read in full and rewritten here.
%% Numbers from tools/figures/s05_numbers.json, keys "blanchard_kahn",
%% "pert_first_order".

%% ECON 282E -- Foundations of Macroeconomics (UC Riverside, Fall 2026)
%% Shared Beamer preamble for every deck in slides/S01 ... slides/S10.
%%
%% Usage (a deck lives in slides/S0X/ and is compiled from that folder):
%%   %% ECON 282E -- Foundations of Macroeconomics (UC Riverside, Fall 2026)
%% Shared Beamer preamble for every deck in slides/S01 ... slides/S10.
%%
%% Usage (a deck lives in slides/S0X/ and is compiled from that folder):
%%   %% ECON 282E -- Foundations of Macroeconomics (UC Riverside, Fall 2026)
%% Shared Beamer preamble for every deck in slides/S01 ... slides/S10.
%%
%% Usage (a deck lives in slides/S0X/ and is compiled from that folder):
%%   \input{../preamble/course_preamble.tex}
%%   \decktitle{1}{A1}{Who I Am \& Why This Course}{September 24, 2026}
%%   \begin{document}
%%   \makedecktitle
%%   ...
%%
%% Adapted from Courses/AI-ECON-2026/source/shared_preamble.tex (Metropolis theme,
%% concept/caution/example/takeaway boxes, \sectiondivider). Changes for this course:
%%   * course title block (\decktitle / \makedecktitle) instead of \makebeamertitle
%%   * \zh{...} typesets Chinese (book titles) under pdflatex via CJKutf8
%%   * researchbox for the "Research in the wild" frames
%%   * section pages off; TikZ libraries and the course map loaded here
%% Build with pdflatex only (two passes); tools/build_slides.py does this for you.

\documentclass[10pt,english,aspectratio=169]{beamer}

\usepackage[T1]{fontenc}
\usepackage{textcomp}
\usepackage[utf8]{inputenc}
\usepackage{babel}
\usepackage{CJKutf8}

\usepackage{amsmath, amssymb, amsthm, graphicx}
\usepackage{booktabs, multirow, tabularx}
\usepackage{tikz}
\usetikzlibrary{positioning,arrows.meta,shapes,calc,fit,backgrounds}
\usepackage{pgfplots}
\pgfplotsset{compat=1.16}
\usepackage[most]{tcolorbox}
\tcbuselibrary{skins, breakable}
\usepackage{listings}
\usepackage{xcolor}

\lstset{
  basicstyle=\ttfamily\small,
  keywordstyle=\color{blue!80!black}\bfseries,
  commentstyle=\color{green!50!black}\itshape,
  stringstyle=\color{orange!80!black},
  backgroundcolor=\color{gray!8},
  frame=single,
  rulecolor=\color{gray!40},
  breaklines=true,
  showstringspaces=false,
  numbers=none,
}

\usetheme{metropolis}
\usefonttheme{professionalfonts}
\metroset{sectionpage=none}

\definecolor{LightGray}{RGB}{230,230,230}
\definecolor{RedTitle}{RGB}{0,0,200}
\definecolor{VeryLightGray}{RGB}{245,245,245}
\definecolor{DarkText}{RGB}{0,0,0}
\definecolor{UCRBlue}{RGB}{0,45,106}
\definecolor{UCRGold}{RGB}{241,171,0}

% Stage colours of the course arc (used by the course map and elsewhere)
\colorlet{StageTrunk}{blue!12}
\colorlet{StageBranch}{cyan!14}
\colorlet{StageMachine}{gray!18}
\colorlet{StageWall}{red!14}
\colorlet{StageTools}{orange!20}
\colorlet{StageData}{green!16}
\colorlet{StageAgents}{violet!14}

\hypersetup{bookmarks=false,breaklinks=true,pdfborder={0 0 0},colorlinks=true,
  linkcolor=DarkText,urlcolor=RedTitle}

\setbeamercolor{background canvas}{bg=white}
\setbeamercolor{title}{fg=RedTitle}
\setbeamercolor{frametitle}{bg=VeryLightGray,fg=DarkText}
\setbeamercolor{normal text}{fg=DarkText}
\setbeamercolor{alerted text}{fg=RedTitle}
\setbeamersize{text margin left=5mm, text margin right=5mm}
\addtobeamertemplate{frametitle}{}{\vspace{-0.5em}}

\setbeamertemplate{navigation symbols}{}
\setbeamertemplate{footline}{%
  \hspace{5mm}{\usebeamerfont{page number in head/foot}\color{black!45}%
  ECON 282E \,$\cdot$\, \insertshortdeck}\hfill%
  \usebeamercolor[fg]{page number in head/foot}%
  \usebeamerfont{page number in head/foot}%
  \insertframenumber\,/\,\inserttotalframenumber\kern1em\vskip2pt%
}

% ---------------------------------------------------------------------------
% Chinese text under pdflatex (book titles etc.)
\newcommand{\zh}[1]{\begin{CJK*}{UTF8}{gbsn}#1\end{CJK*}}

% ---------------------------------------------------------------------------
% Deck title block
%   \decktitle{session}{deck id}{title}{date}
\newcommand{\insertshortdeck}{}
\newcommand{\decksession}{}
\newcommand{\deckid}{}
\newcommand{\decktitletext}{}
\newcommand{\deckdate}{}
\newcommand{\decktitle}[4]{%
  \renewcommand{\decksession}{#1}%
  \renewcommand{\deckid}{#2}%
  \renewcommand{\decktitletext}{#3}%
  \renewcommand{\deckdate}{#4}%
  \renewcommand{\insertshortdeck}{Session #1 \,$\cdot$\, Deck #1.#2}%
  \title{#3}%
  \author{Zhigang Feng}%
  \date{#4}%
}
\newcommand{\makedecktitle}{%
  \begin{frame}[plain,noframenumbering]
    \begin{tikzpicture}[remember picture,overlay]
      \fill[UCRBlue] (current page.north west) rectangle ([yshift=-0.55cm]current page.north east);
      \fill[UCRGold] ([yshift=-0.55cm]current page.north west) rectangle ([yshift=-0.62cm]current page.north east);
      \node[anchor=west, text=white, font=\small\bfseries] at ([xshift=5mm,yshift=-0.275cm]current page.north west)
        {ECON 282E \,$\cdot$\, Foundations of Macroeconomics};
      \node[anchor=east, text=white, font=\small] at ([xshift=-5mm,yshift=-0.275cm]current page.north east)
        {UC Riverside \,$\cdot$\, Fall 2026};
    \end{tikzpicture}
    \vspace*{1.1cm}
    {\usebeamerfont{title}\color{black!55}\large Session \decksession\ \,$\cdot$\, Deck \decksession.\deckid\par}
    \vspace{0.25cm}
    {\usebeamerfont{title}\color{RedTitle}\LARGE\bfseries \decktitletext\par}
    \vspace{0.7cm}
    {\large Zhigang Feng\par}
    \vspace{0.15cm}
    {\color{black!65}\deckdate\par}
    \vfill
    {\footnotesize\itshape\color{black!55}Build it on paper $\;\to\;$ solve it on the machine $\;\to\;$ push it to the frontier with AI\par}
  \end{frame}}

% ---------------------------------------------------------------------------
% Section divider (plain frame, not counted)
\newcommand{\sectiondivider}[2]{%
\begin{frame}[plain,noframenumbering]
\begin{center}
\vfill
{\Large\textbf{#1}\par}
\vspace{0.35cm}
{\normalsize #2\par}
\vfill
\end{center}
\end{frame}
}

% ---------------------------------------------------------------------------
% Boxes
\newtcolorbox{conceptbox}[1][]{colback=blue!3, colframe=RedTitle!55, boxrule=0.35mm,
  arc=1mm, left=2mm, right=2mm, top=1mm, bottom=1mm, #1}
\newtcolorbox{cautionbox}[1][]{colback=red!3, colframe=red!50!black, boxrule=0.35mm,
  arc=1mm, left=2mm, right=2mm, top=1mm, bottom=1mm, #1}
\newtcolorbox{examplebox}[1][]{colback=green!3, colframe=green!45!black, boxrule=0.35mm,
  arc=1mm, left=2mm, right=2mm, top=1mm, bottom=1mm, #1}
\newtcolorbox{takeawaybox}[1][]{colback=VeryLightGray, colframe=RedTitle!55, boxrule=0.35mm,
  arc=1mm, left=2mm, right=2mm, top=1mm, bottom=1mm, #1}
\newtcolorbox{researchbox}[1][]{enhanced, colback=violet!4, colframe=violet!60!black,
  boxrule=0.35mm, arc=1mm, left=2mm, right=2mm, top=1mm, bottom=1mm,
  fonttitle=\bfseries\small, title={Research in the wild}, #1}

\newcommand{\compacttables}{\renewcommand{\arraystretch}{1.15}}
\newcommand{\src}[1]{{\tiny\color{black!55}#1}}

% ---------------------------------------------------------------------------
\input{../preamble/course_map.tex}

%%   \decktitle{1}{A1}{Who I Am \& Why This Course}{September 24, 2026}
%%   \begin{document}
%%   \makedecktitle
%%   ...
%%
%% Adapted from Courses/AI-ECON-2026/source/shared_preamble.tex (Metropolis theme,
%% concept/caution/example/takeaway boxes, \sectiondivider). Changes for this course:
%%   * course title block (\decktitle / \makedecktitle) instead of \makebeamertitle
%%   * \zh{...} typesets Chinese (book titles) under pdflatex via CJKutf8
%%   * researchbox for the "Research in the wild" frames
%%   * section pages off; TikZ libraries and the course map loaded here
%% Build with pdflatex only (two passes); tools/build_slides.py does this for you.

\documentclass[10pt,english,aspectratio=169]{beamer}

\usepackage[T1]{fontenc}
\usepackage{textcomp}
\usepackage[utf8]{inputenc}
\usepackage{babel}
\usepackage{CJKutf8}

\usepackage{amsmath, amssymb, amsthm, graphicx}
\usepackage{booktabs, multirow, tabularx}
\usepackage{tikz}
\usetikzlibrary{positioning,arrows.meta,shapes,calc,fit,backgrounds}
\usepackage{pgfplots}
\pgfplotsset{compat=1.16}
\usepackage[most]{tcolorbox}
\tcbuselibrary{skins, breakable}
\usepackage{listings}
\usepackage{xcolor}

\lstset{
  basicstyle=\ttfamily\small,
  keywordstyle=\color{blue!80!black}\bfseries,
  commentstyle=\color{green!50!black}\itshape,
  stringstyle=\color{orange!80!black},
  backgroundcolor=\color{gray!8},
  frame=single,
  rulecolor=\color{gray!40},
  breaklines=true,
  showstringspaces=false,
  numbers=none,
}

\usetheme{metropolis}
\usefonttheme{professionalfonts}
\metroset{sectionpage=none}

\definecolor{LightGray}{RGB}{230,230,230}
\definecolor{RedTitle}{RGB}{0,0,200}
\definecolor{VeryLightGray}{RGB}{245,245,245}
\definecolor{DarkText}{RGB}{0,0,0}
\definecolor{UCRBlue}{RGB}{0,45,106}
\definecolor{UCRGold}{RGB}{241,171,0}

% Stage colours of the course arc (used by the course map and elsewhere)
\colorlet{StageTrunk}{blue!12}
\colorlet{StageBranch}{cyan!14}
\colorlet{StageMachine}{gray!18}
\colorlet{StageWall}{red!14}
\colorlet{StageTools}{orange!20}
\colorlet{StageData}{green!16}
\colorlet{StageAgents}{violet!14}

\hypersetup{bookmarks=false,breaklinks=true,pdfborder={0 0 0},colorlinks=true,
  linkcolor=DarkText,urlcolor=RedTitle}

\setbeamercolor{background canvas}{bg=white}
\setbeamercolor{title}{fg=RedTitle}
\setbeamercolor{frametitle}{bg=VeryLightGray,fg=DarkText}
\setbeamercolor{normal text}{fg=DarkText}
\setbeamercolor{alerted text}{fg=RedTitle}
\setbeamersize{text margin left=5mm, text margin right=5mm}
\addtobeamertemplate{frametitle}{}{\vspace{-0.5em}}

\setbeamertemplate{navigation symbols}{}
\setbeamertemplate{footline}{%
  \hspace{5mm}{\usebeamerfont{page number in head/foot}\color{black!45}%
  ECON 282E \,$\cdot$\, \insertshortdeck}\hfill%
  \usebeamercolor[fg]{page number in head/foot}%
  \usebeamerfont{page number in head/foot}%
  \insertframenumber\,/\,\inserttotalframenumber\kern1em\vskip2pt%
}

% ---------------------------------------------------------------------------
% Chinese text under pdflatex (book titles etc.)
\newcommand{\zh}[1]{\begin{CJK*}{UTF8}{gbsn}#1\end{CJK*}}

% ---------------------------------------------------------------------------
% Deck title block
%   \decktitle{session}{deck id}{title}{date}
\newcommand{\insertshortdeck}{}
\newcommand{\decksession}{}
\newcommand{\deckid}{}
\newcommand{\decktitletext}{}
\newcommand{\deckdate}{}
\newcommand{\decktitle}[4]{%
  \renewcommand{\decksession}{#1}%
  \renewcommand{\deckid}{#2}%
  \renewcommand{\decktitletext}{#3}%
  \renewcommand{\deckdate}{#4}%
  \renewcommand{\insertshortdeck}{Session #1 \,$\cdot$\, Deck #1.#2}%
  \title{#3}%
  \author{Zhigang Feng}%
  \date{#4}%
}
\newcommand{\makedecktitle}{%
  \begin{frame}[plain,noframenumbering]
    \begin{tikzpicture}[remember picture,overlay]
      \fill[UCRBlue] (current page.north west) rectangle ([yshift=-0.55cm]current page.north east);
      \fill[UCRGold] ([yshift=-0.55cm]current page.north west) rectangle ([yshift=-0.62cm]current page.north east);
      \node[anchor=west, text=white, font=\small\bfseries] at ([xshift=5mm,yshift=-0.275cm]current page.north west)
        {ECON 282E \,$\cdot$\, Foundations of Macroeconomics};
      \node[anchor=east, text=white, font=\small] at ([xshift=-5mm,yshift=-0.275cm]current page.north east)
        {UC Riverside \,$\cdot$\, Fall 2026};
    \end{tikzpicture}
    \vspace*{1.1cm}
    {\usebeamerfont{title}\color{black!55}\large Session \decksession\ \,$\cdot$\, Deck \decksession.\deckid\par}
    \vspace{0.25cm}
    {\usebeamerfont{title}\color{RedTitle}\LARGE\bfseries \decktitletext\par}
    \vspace{0.7cm}
    {\large Zhigang Feng\par}
    \vspace{0.15cm}
    {\color{black!65}\deckdate\par}
    \vfill
    {\footnotesize\itshape\color{black!55}Build it on paper $\;\to\;$ solve it on the machine $\;\to\;$ push it to the frontier with AI\par}
  \end{frame}}

% ---------------------------------------------------------------------------
% Section divider (plain frame, not counted)
\newcommand{\sectiondivider}[2]{%
\begin{frame}[plain,noframenumbering]
\begin{center}
\vfill
{\Large\textbf{#1}\par}
\vspace{0.35cm}
{\normalsize #2\par}
\vfill
\end{center}
\end{frame}
}

% ---------------------------------------------------------------------------
% Boxes
\newtcolorbox{conceptbox}[1][]{colback=blue!3, colframe=RedTitle!55, boxrule=0.35mm,
  arc=1mm, left=2mm, right=2mm, top=1mm, bottom=1mm, #1}
\newtcolorbox{cautionbox}[1][]{colback=red!3, colframe=red!50!black, boxrule=0.35mm,
  arc=1mm, left=2mm, right=2mm, top=1mm, bottom=1mm, #1}
\newtcolorbox{examplebox}[1][]{colback=green!3, colframe=green!45!black, boxrule=0.35mm,
  arc=1mm, left=2mm, right=2mm, top=1mm, bottom=1mm, #1}
\newtcolorbox{takeawaybox}[1][]{colback=VeryLightGray, colframe=RedTitle!55, boxrule=0.35mm,
  arc=1mm, left=2mm, right=2mm, top=1mm, bottom=1mm, #1}
\newtcolorbox{researchbox}[1][]{enhanced, colback=violet!4, colframe=violet!60!black,
  boxrule=0.35mm, arc=1mm, left=2mm, right=2mm, top=1mm, bottom=1mm,
  fonttitle=\bfseries\small, title={Research in the wild}, #1}

\newcommand{\compacttables}{\renewcommand{\arraystretch}{1.15}}
\newcommand{\src}[1]{{\tiny\color{black!55}#1}}

% ---------------------------------------------------------------------------
%% Course map: the recurring "you are here" slide for ECON 282E.
%% Loaded by course_preamble.tex. Pattern follows AI-ECON-2026 lec01_map.tex, but the map
%% shows the whole ten-session course rather than one lecture.
%%
%%   \CourseMap{s}{label}   s = 1..10 highlights session s and prints "you are here: label"
%%                          (e.g. \CourseMap{1}{1.A2}); earlier sessions get a check mark.
%%   \CourseMap{0}{}        full overview, nothing highlighted.
%%
%% Note: TikZ option lists are comma-split before TeX conditionals run, so every \ifnum
%% below sits outside [...] option lists.

\newcommand{\CourseMap}[2]{%
\begin{center}
\begin{tikzpicture}[
    sess/.style={rectangle, rounded corners=2pt, draw=black!45, minimum width=1.34cm,
        minimum height=1.12cm, text width=1.26cm, align=center, inner sep=1pt},
    stage/.style={font=\tiny\bfseries, text=black!70},
]
  \foreach \i/\d/\t/\c in {%
      1/Sep 24/Intro \\ Growth/StageTrunk,
      2/Oct 1/HA $\cdot$ OLG \\ Policy/StageBranch,
      3/Oct 8/Num.\ DP \\ Python/StageMachine,
      4/Oct 15/Numerics \\ I/StageMachine,
      5/Oct 22/Numerics \\ II/StageMachine,
      6/Oct 29/The wall \\ \textbf{Midterm}/StageWall,
      7/Nov 5/ML \\ Deep nets/StageTools,
      8/Nov 12/RL \\ HA $+$ DL/StageTools,
      9/Nov 19/Text \\ LLMs/StageData,
      10/Dec 3/Agents \\ \textbf{Projects}/StageAgents} {
    \node[sess, fill=\c] (n\i) at ({1.46*(\i-1)}, 0)
      {{\scriptsize\bfseries S\i}\\[-1pt]{\tiny\color{black!60}\d}\\[1pt]{\tiny \t}};
    \ifnum\i<#1
      \node[font=\scriptsize, text=green!45!black] at ([xshift=-2.5mm,yshift=-2.2mm]n\i.north east) {$\checkmark$};
    \fi
    \ifnum\i=#1
      \draw[RedTitle, very thick, rounded corners=2pt]
        ([xshift=-0.6mm,yshift=-0.6mm]n\i.south west) rectangle ([xshift=0.6mm,yshift=0.6mm]n\i.north east);
      % keep the label inside the picture at both ends of the row
      \ifnum\i<3
        \node[font=\scriptsize\bfseries, text=RedTitle, anchor=south west] at ([yshift=1.2mm]n\i.north west)
          {$\blacktriangledown$\,you are here: #2};
      \else\ifnum\i>8
        \node[font=\scriptsize\bfseries, text=RedTitle, anchor=south east] at ([yshift=1.2mm]n\i.north east)
          {you are here: #2\,$\blacktriangledown$};
      \else
        \node[font=\scriptsize\bfseries, text=RedTitle, anchor=south] at ([yshift=1.2mm]n\i.north)
          {you are here: #2\,$\blacktriangledown$};
      \fi\fi
    \fi
  }
  % stage band under the sessions
  \foreach \a/\b/\lab/\c in {1/1/trunk/StageTrunk, 2/2/branches/StageBranch,
      3/5/the machine $\cdot$ classical methods/StageMachine, 6/6/the wall/StageWall,
      7/8/new tools/StageTools, 9/9/new data/StageData, 10/10/colleagues/StageAgents} {
    \fill[\c] ([yshift=-1.5mm]n\a.south west) rectangle ([yshift=-4.6mm]n\b.south east);
    \path ([yshift=-1.5mm]n\a.south west) -- ([yshift=-4.6mm]n\b.south east)
      node[midway, stage] {\lab};
  }
  \node[font=\footnotesize\itshape, text=black!70] at ([yshift=-9.5mm]$(n1.south)!0.5!(n10.south)$)
    {Build it on paper $\;\to\;$ solve it on the machine $\;\to\;$ push it to the frontier with AI};
\end{tikzpicture}
\end{center}
}


%%   \decktitle{1}{A1}{Who I Am \& Why This Course}{September 24, 2026}
%%   \begin{document}
%%   \makedecktitle
%%   ...
%%
%% Adapted from Courses/AI-ECON-2026/source/shared_preamble.tex (Metropolis theme,
%% concept/caution/example/takeaway boxes, \sectiondivider). Changes for this course:
%%   * course title block (\decktitle / \makedecktitle) instead of \makebeamertitle
%%   * \zh{...} typesets Chinese (book titles) under pdflatex via CJKutf8
%%   * researchbox for the "Research in the wild" frames
%%   * section pages off; TikZ libraries and the course map loaded here
%% Build with pdflatex only (two passes); tools/build_slides.py does this for you.

\documentclass[10pt,english,aspectratio=169]{beamer}

\usepackage[T1]{fontenc}
\usepackage{textcomp}
\usepackage[utf8]{inputenc}
\usepackage{babel}
\usepackage{CJKutf8}

\usepackage{amsmath, amssymb, amsthm, graphicx}
\usepackage{booktabs, multirow, tabularx}
\usepackage{tikz}
\usetikzlibrary{positioning,arrows.meta,shapes,calc,fit,backgrounds}
\usepackage{pgfplots}
\pgfplotsset{compat=1.16}
\usepackage[most]{tcolorbox}
\tcbuselibrary{skins, breakable}
\usepackage{listings}
\usepackage{xcolor}

\lstset{
  basicstyle=\ttfamily\small,
  keywordstyle=\color{blue!80!black}\bfseries,
  commentstyle=\color{green!50!black}\itshape,
  stringstyle=\color{orange!80!black},
  backgroundcolor=\color{gray!8},
  frame=single,
  rulecolor=\color{gray!40},
  breaklines=true,
  showstringspaces=false,
  numbers=none,
}

\usetheme{metropolis}
\usefonttheme{professionalfonts}
\metroset{sectionpage=none}

\definecolor{LightGray}{RGB}{230,230,230}
\definecolor{RedTitle}{RGB}{0,0,200}
\definecolor{VeryLightGray}{RGB}{245,245,245}
\definecolor{DarkText}{RGB}{0,0,0}
\definecolor{UCRBlue}{RGB}{0,45,106}
\definecolor{UCRGold}{RGB}{241,171,0}

% Stage colours of the course arc (used by the course map and elsewhere)
\colorlet{StageTrunk}{blue!12}
\colorlet{StageBranch}{cyan!14}
\colorlet{StageMachine}{gray!18}
\colorlet{StageWall}{red!14}
\colorlet{StageTools}{orange!20}
\colorlet{StageData}{green!16}
\colorlet{StageAgents}{violet!14}

\hypersetup{bookmarks=false,breaklinks=true,pdfborder={0 0 0},colorlinks=true,
  linkcolor=DarkText,urlcolor=RedTitle}

\setbeamercolor{background canvas}{bg=white}
\setbeamercolor{title}{fg=RedTitle}
\setbeamercolor{frametitle}{bg=VeryLightGray,fg=DarkText}
\setbeamercolor{normal text}{fg=DarkText}
\setbeamercolor{alerted text}{fg=RedTitle}
\setbeamersize{text margin left=5mm, text margin right=5mm}
\addtobeamertemplate{frametitle}{}{\vspace{-0.5em}}

\setbeamertemplate{navigation symbols}{}
\setbeamertemplate{footline}{%
  \hspace{5mm}{\usebeamerfont{page number in head/foot}\color{black!45}%
  ECON 282E \,$\cdot$\, \insertshortdeck}\hfill%
  \usebeamercolor[fg]{page number in head/foot}%
  \usebeamerfont{page number in head/foot}%
  \insertframenumber\,/\,\inserttotalframenumber\kern1em\vskip2pt%
}

% ---------------------------------------------------------------------------
% Chinese text under pdflatex (book titles etc.)
\newcommand{\zh}[1]{\begin{CJK*}{UTF8}{gbsn}#1\end{CJK*}}

% ---------------------------------------------------------------------------
% Deck title block
%   \decktitle{session}{deck id}{title}{date}
\newcommand{\insertshortdeck}{}
\newcommand{\decksession}{}
\newcommand{\deckid}{}
\newcommand{\decktitletext}{}
\newcommand{\deckdate}{}
\newcommand{\decktitle}[4]{%
  \renewcommand{\decksession}{#1}%
  \renewcommand{\deckid}{#2}%
  \renewcommand{\decktitletext}{#3}%
  \renewcommand{\deckdate}{#4}%
  \renewcommand{\insertshortdeck}{Session #1 \,$\cdot$\, Deck #1.#2}%
  \title{#3}%
  \author{Zhigang Feng}%
  \date{#4}%
}
\newcommand{\makedecktitle}{%
  \begin{frame}[plain,noframenumbering]
    \begin{tikzpicture}[remember picture,overlay]
      \fill[UCRBlue] (current page.north west) rectangle ([yshift=-0.55cm]current page.north east);
      \fill[UCRGold] ([yshift=-0.55cm]current page.north west) rectangle ([yshift=-0.62cm]current page.north east);
      \node[anchor=west, text=white, font=\small\bfseries] at ([xshift=5mm,yshift=-0.275cm]current page.north west)
        {ECON 282E \,$\cdot$\, Foundations of Macroeconomics};
      \node[anchor=east, text=white, font=\small] at ([xshift=-5mm,yshift=-0.275cm]current page.north east)
        {UC Riverside \,$\cdot$\, Fall 2026};
    \end{tikzpicture}
    \vspace*{1.1cm}
    {\usebeamerfont{title}\color{black!55}\large Session \decksession\ \,$\cdot$\, Deck \decksession.\deckid\par}
    \vspace{0.25cm}
    {\usebeamerfont{title}\color{RedTitle}\LARGE\bfseries \decktitletext\par}
    \vspace{0.7cm}
    {\large Zhigang Feng\par}
    \vspace{0.15cm}
    {\color{black!65}\deckdate\par}
    \vfill
    {\footnotesize\itshape\color{black!55}Build it on paper $\;\to\;$ solve it on the machine $\;\to\;$ push it to the frontier with AI\par}
  \end{frame}}

% ---------------------------------------------------------------------------
% Section divider (plain frame, not counted)
\newcommand{\sectiondivider}[2]{%
\begin{frame}[plain,noframenumbering]
\begin{center}
\vfill
{\Large\textbf{#1}\par}
\vspace{0.35cm}
{\normalsize #2\par}
\vfill
\end{center}
\end{frame}
}

% ---------------------------------------------------------------------------
% Boxes
\newtcolorbox{conceptbox}[1][]{colback=blue!3, colframe=RedTitle!55, boxrule=0.35mm,
  arc=1mm, left=2mm, right=2mm, top=1mm, bottom=1mm, #1}
\newtcolorbox{cautionbox}[1][]{colback=red!3, colframe=red!50!black, boxrule=0.35mm,
  arc=1mm, left=2mm, right=2mm, top=1mm, bottom=1mm, #1}
\newtcolorbox{examplebox}[1][]{colback=green!3, colframe=green!45!black, boxrule=0.35mm,
  arc=1mm, left=2mm, right=2mm, top=1mm, bottom=1mm, #1}
\newtcolorbox{takeawaybox}[1][]{colback=VeryLightGray, colframe=RedTitle!55, boxrule=0.35mm,
  arc=1mm, left=2mm, right=2mm, top=1mm, bottom=1mm, #1}
\newtcolorbox{researchbox}[1][]{enhanced, colback=violet!4, colframe=violet!60!black,
  boxrule=0.35mm, arc=1mm, left=2mm, right=2mm, top=1mm, bottom=1mm,
  fonttitle=\bfseries\small, title={Research in the wild}, #1}

\newcommand{\compacttables}{\renewcommand{\arraystretch}{1.15}}
\newcommand{\src}[1]{{\tiny\color{black!55}#1}}

% ---------------------------------------------------------------------------
%% Course map: the recurring "you are here" slide for ECON 282E.
%% Loaded by course_preamble.tex. Pattern follows AI-ECON-2026 lec01_map.tex, but the map
%% shows the whole ten-session course rather than one lecture.
%%
%%   \CourseMap{s}{label}   s = 1..10 highlights session s and prints "you are here: label"
%%                          (e.g. \CourseMap{1}{1.A2}); earlier sessions get a check mark.
%%   \CourseMap{0}{}        full overview, nothing highlighted.
%%
%% Note: TikZ option lists are comma-split before TeX conditionals run, so every \ifnum
%% below sits outside [...] option lists.

\newcommand{\CourseMap}[2]{%
\begin{center}
\begin{tikzpicture}[
    sess/.style={rectangle, rounded corners=2pt, draw=black!45, minimum width=1.34cm,
        minimum height=1.12cm, text width=1.26cm, align=center, inner sep=1pt},
    stage/.style={font=\tiny\bfseries, text=black!70},
]
  \foreach \i/\d/\t/\c in {%
      1/Sep 24/Intro \\ Growth/StageTrunk,
      2/Oct 1/HA $\cdot$ OLG \\ Policy/StageBranch,
      3/Oct 8/Num.\ DP \\ Python/StageMachine,
      4/Oct 15/Numerics \\ I/StageMachine,
      5/Oct 22/Numerics \\ II/StageMachine,
      6/Oct 29/The wall \\ \textbf{Midterm}/StageWall,
      7/Nov 5/ML \\ Deep nets/StageTools,
      8/Nov 12/RL \\ HA $+$ DL/StageTools,
      9/Nov 19/Text \\ LLMs/StageData,
      10/Dec 3/Agents \\ \textbf{Projects}/StageAgents} {
    \node[sess, fill=\c] (n\i) at ({1.46*(\i-1)}, 0)
      {{\scriptsize\bfseries S\i}\\[-1pt]{\tiny\color{black!60}\d}\\[1pt]{\tiny \t}};
    \ifnum\i<#1
      \node[font=\scriptsize, text=green!45!black] at ([xshift=-2.5mm,yshift=-2.2mm]n\i.north east) {$\checkmark$};
    \fi
    \ifnum\i=#1
      \draw[RedTitle, very thick, rounded corners=2pt]
        ([xshift=-0.6mm,yshift=-0.6mm]n\i.south west) rectangle ([xshift=0.6mm,yshift=0.6mm]n\i.north east);
      % keep the label inside the picture at both ends of the row
      \ifnum\i<3
        \node[font=\scriptsize\bfseries, text=RedTitle, anchor=south west] at ([yshift=1.2mm]n\i.north west)
          {$\blacktriangledown$\,you are here: #2};
      \else\ifnum\i>8
        \node[font=\scriptsize\bfseries, text=RedTitle, anchor=south east] at ([yshift=1.2mm]n\i.north east)
          {you are here: #2\,$\blacktriangledown$};
      \else
        \node[font=\scriptsize\bfseries, text=RedTitle, anchor=south] at ([yshift=1.2mm]n\i.north)
          {you are here: #2\,$\blacktriangledown$};
      \fi\fi
    \fi
  }
  % stage band under the sessions
  \foreach \a/\b/\lab/\c in {1/1/trunk/StageTrunk, 2/2/branches/StageBranch,
      3/5/the machine $\cdot$ classical methods/StageMachine, 6/6/the wall/StageWall,
      7/8/new tools/StageTools, 9/9/new data/StageData, 10/10/colleagues/StageAgents} {
    \fill[\c] ([yshift=-1.5mm]n\a.south west) rectangle ([yshift=-4.6mm]n\b.south east);
    \path ([yshift=-1.5mm]n\a.south west) -- ([yshift=-4.6mm]n\b.south east)
      node[midway, stage] {\lab};
  }
  \node[font=\footnotesize\itshape, text=black!70] at ([yshift=-9.5mm]$(n1.south)!0.5!(n10.south)$)
    {Build it on paper $\;\to\;$ solve it on the machine $\;\to\;$ push it to the frontier with AI};
\end{tikzpicture}
\end{center}
}


\graphicspath{{./figures/}}

\lstset{language=Python, basicstyle=\ttfamily\scriptsize, frame=none,
        backgroundcolor=\color{gray!8}, aboveskip=2pt, belowskip=2pt}

\decktitle{5}{A4}{Linear Rational-Expectations Systems}{Thursday, October 22, 2026}

\begin{document}

\makedecktitle

%=============================================================================
\begin{frame}{Where We Are}
\CourseMap{5}{5.A4}
\vspace{-1mm}
\begin{center}
\small \textbf{Block A:} 5.A1 what perturbation computes $\;\cdot\;$ 5.A2 the model and the parameter $\;\cdot\;$ 5.A3 first order $\;\cdot\;$ \textbf{5.A4} linear RE systems $\;\cdot\;$ 5.A5 higher order, risk, pruning $\;\cdot\;$ 5.A6 the worked example.
\end{center}
\end{frame}

%=============================================================================
\begin{frame}{Linearize, or Log-Linearize?}
\begin{columns}[T]
\begin{column}{0.55\textwidth}
\small
Both are first-order Taylor expansions; they differ in the variable.
\medskip

\textbf{Log-linearization.} Substitute $x_t = x\,e^{\hat{x}_t}$ where
\[
  \hat{x}_t = \ln\frac{x_t}{x} ,
\]
and expand in $\hat{x}_t$. Then every coefficient is an \textbf{elasticity} and every deviation is a \textbf{percentage}, which is how macro results are reported anyway.
\medskip

\textbf{Linearization.} Expand in $x_t-x$ directly. Coefficients carry units.
\medskip

They are the same approximation under a change of variables, so you may do either and convert afterwards --- but the accuracy is \emph{not} identical, because the two are exact for different functions.
\end{column}
\begin{column}{0.42\textwidth}
\begin{examplebox}[title=\textbf{Which is better here?}]\small
5.A2 settled it for this model: the Brock--Mirman policy is \textbf{exactly linear in logs} and curved in levels. Log-linearizing recovers it exactly; linearizing does not.
\end{examplebox}
\medskip
\begin{cautionbox}\footnotesize
That is a property of this model, not a law. Logs are the right default in macro because most variables are positive and grow.
\end{cautionbox}
\end{column}
\end{columns}
\vspace{1mm}
\src{Fern\'andez-Villaverde \& Guerr\'on-Quintana, \emph{Solution Methods} V (Appendix: Linearization), ``Linearization I''.}
\end{frame}

%=============================================================================
\begin{frame}{Linearizing the Equilibrium Conditions}
\begin{columns}[T]
\begin{column}{0.55\textwidth}
\small
Take the quarterly model of 5.A3 and expand each condition around $(k^{*},c^{*},0)$, writing hats for log deviations.
\medskip

\textbf{Resource constraint.}
\[
  k^{*}\hat{k}_{t+1}
  = \big(\alpha y^{*}+(1-\delta)k^{*}\big)\hat{k}_{t} + y^{*}z_t - c^{*}\hat{c}_t
\]
\textbf{Euler equation.}
\[
  \mathbb{E}_t\hat{c}_{t+1}
  = \hat{c}_t + \kappa\,\mathbb{E}_t z_{t+1} + \kappa(\alpha-1)\hat{k}_{t+1}
\]
with $\kappa=\beta\alpha (k^{*})^{\alpha-1}$ --- the share of the gross return that actually moves. Measured: $\kappa=0.03475$.
\medskip

\textbf{Shock.} $\;\mathbb{E}_t z_{t+1}=\rho z_t$.
\end{column}
\begin{column}{0.42\textwidth}
\begin{conceptbox}[title=\textbf{Why $\kappa$ is small}]\small
At $\delta=0.025$, most of next period's gross return is undepreciated capital, which does not respond to $z$ or $k$. Only $3.5\%$ of it does.
\medskip
That single number is why capital is so persistent here, and why the model propagates shocks for years.
\end{conceptbox}
\end{column}
\end{columns}
\vspace{1mm}
\src{Structure follows FV\&G, \emph{Solution Methods} V (Appendix), ``Linearization II--III''; instantiated on this course's model rather than theirs, which carries endogenous labour. Ledger \texttt{pert\_first\_order}.}
\end{frame}

%=============================================================================
\begin{frame}{Rewriting It as a System}
\begin{columns}[T]
\begin{column}{0.55\textwidth}
\small
After collecting terms, any linearized model has the shape
\begin{align*}
 A\,\hat{s}_{t+1} + B\,\hat{s}_{t} + C\,\hat{y}_{t} + D z_t &= 0\\
 \mathbb{E}_t\big(F\hat{s}_{t+1}+G\hat{s}_t+H\hat{s}_{t-1}\\
 \qquad +J\hat{y}_{t+1}+K\hat{y}_t+Lz_{t+1}+Mz_t\big)&=0\\
 \mathbb{E}_t z_{t+1} &= N z_t
\end{align*}
with $m$ states $\hat{s}_t$, $n$ controls $\hat{y}_t$, and $k$ exogenous processes $z_t$.
\medskip

The first block is the conditions that hold \textbf{within} a period; the second is the ones involving \textbf{expectations}. That split is the whole difficulty: an ordinary difference equation would already be solved.
\end{column}
\begin{column}{0.42\textwidth}
\begin{examplebox}\footnotesize
You may eliminate controls by substitution before forming the system, as FV\&G do with $c_t$. It is not necessary, and it changes only the arithmetic --- not the policy functions you end up with.
\end{examplebox}
\medskip
\begin{cautionbox}\footnotesize
$N$ must have only stable eigenvalues: the exogenous process has to be stationary, or nothing below works.
\end{cautionbox}
\end{column}
\end{columns}
\vspace{1mm}
\src{FV\&G, \emph{Solution Methods} V (Appendix), ``General structure of linearized system''.}
\end{frame}

%=============================================================================
\begin{frame}{Undetermined Coefficients: Guess, Substitute, Match}
\begin{columns}[T]
\begin{column}{0.55\textwidth}
\small
Guess that the solution is linear:
\begin{align*}
 \hat{s}_t &= P\,\hat{s}_{t-1} + Q\,z_t\\
 \hat{y}_t &= R\,\hat{s}_{t-1} + U\,z_t
\end{align*}
Substitute into the system, use linearity of expectations and $\mathbb{E}_t z_{t+1}=Nz_t$, and collect.
\medskip

Because the result must hold for \textbf{every} value of $\hat{s}_{t-1}$ and $z_t$, each coefficient must vanish separately. In the scalar case that gives, on $\hat{s}_{t-1}$,
\[
  AP_1+B+CR_1=0,\qquad GP_1+H+JR_1P_1+KR_1=0,
\]
and two more on $z_t$.
\end{column}
\begin{column}{0.42\textwidth}
\begin{conceptbox}[title=\textbf{The same trick as 5.A3}]\small
``It holds identically, so match coefficients.'' There we matched derivatives of $F$; here we match powers of the state. Same idea, different bookkeeping.
\end{conceptbox}
\medskip
\begin{takeawaybox}\footnotesize
Four equations, four unknowns $(P_1,Q,R_1,U)$ --- and, crucially, the first pair does \emph{not} involve $Q$ or $U$. The state block solves on its own.
\end{takeawaybox}
\end{column}
\end{columns}
\vspace{1mm}
\src{FV\&G, \emph{Solution Methods} V (Appendix), ``Undetermined coefficients'' and ``Solving the system I''.}
\end{frame}

%=============================================================================
\begin{frame}{The Quadratic, and Its Two Roots}
\begin{columns}[T]
\begin{column}{0.55\textwidth}
\small
Eliminating $R_1=-C^{-1}(AP_1+B)$ from the pair leaves a single \textbf{quadratic} equation in $P_1$:
\[
  P_1^{2}+\Big(\tfrac{B}{A}+\tfrac{K}{J}-\tfrac{GC}{JA}\Big)P_1
  +\frac{KB-HC}{JA}=0 .
\]
It has two roots,
\[
  P_1=\tfrac12\Big[-\big(\cdot\big)\pm\sqrt{(\cdot)^2-4\tfrac{KB-HC}{JA}}\Big],
\]
\textbf{one stable and one unstable}. Keep the stable one, recover $R_1$, and $Q$ and $U$ then follow from two linear equations.
\medskip

This is where the ``quadratic system'' of 5.A3 turns into an actual choice --- and why counting equations was never enough.
\end{column}
\begin{column}{0.42\textwidth}
\begin{conceptbox}[title=\textbf{Why quadratic at all?}]\small
Because the model has a \textbf{stable and an unstable manifold}. The two roots are those two manifolds; discarding the explosive one is imposing the transversality condition.
\end{conceptbox}
\medskip
\begin{examplebox}\footnotesize
1.B1's saddle path, in algebra. There we drew it; here we solve for it.
\end{examplebox}
\end{column}
\end{columns}
\vspace{1mm}
\src{FV\&G, \emph{Solution Methods} V (Appendix), ``Solving the system II--III''; and \emph{Solution Methods} V, ``Solving the system II'' (``Why quadratic? Stable and unstable manifold'').}
\end{frame}

%=============================================================================
\begin{frame}{In Matrices: One Quadratic Equation}
\begin{columns}[T]
\begin{column}{0.55\textwidth}
\small
With more than one state the same derivation gives a \textbf{matrix} quadratic. $P$ solves
\begin{align*}
  \big(F-JC^{-1}A\big)P^{2}
  &-\big(JC^{-1}B-G+KC^{-1}A\big)P\\
  &-KC^{-1}B+H \;=\; 0 ,
\end{align*}
and the equilibrium is stable if and only if
\[
  \max\big|\mathrm{eig}(P)\big| < 1 .
\]
The rest is linear algebra:
\[
  R=-C^{-1}(AP+B), \qquad U=-C^{-1}(AQ+D),
\]
with $Q$ from a Sylvester equation in $\mathrm{vec}(Q)$.
\end{column}
\begin{column}{0.42\textwidth}
\begin{takeawaybox}\footnotesize
Everything about the model is now in \textbf{one matrix equation} and \textbf{one eigenvalue condition}. The economics has been fully absorbed into $A,\dots,N$.
\end{takeawaybox}
\medskip
\begin{cautionbox}\footnotesize
$C$ must be invertible (rank $n$). When it is not --- more static conditions than controls --- you need the general case, which is the QZ route on the next slide.
\end{cautionbox}
\end{column}
\end{columns}
\vspace{1mm}
\src{FV\&G, \emph{Solution Methods} V (Appendix), ``Policy functions II''. Uhlig's toolkit covers the general $l>n$ case.}
\end{frame}

%=============================================================================
\begin{frame}{How to Solve a Matrix Quadratic}
\begin{columns}[T]
\begin{column}{0.56\textwidth}
\small
To solve $\Psi P^{2}-\Gamma P-\Theta=0$ for the $m\times m$ matrix $P$:
\medskip
\begin{enumerate}\small
  \item Build the $2m\times 2m$ pair
  \[
    \Xi=\begin{bmatrix}\Gamma & \Theta\\ I_m & 0_m\end{bmatrix},
    \qquad
    \Delta=\begin{bmatrix}\Psi & 0_m\\ 0_m & I_m\end{bmatrix}.
  \]
  \item Take the \textbf{generalized} eigenvalues $\omega_i$ and eigenvectors $s_i$ of $\Xi$ with respect to $\Delta$. Each has the form $s_i'=[\omega_i x_i',\,x_i']$.
  \item If $m$ of them give linearly independent $x_i$, then
  \[
    P=\Omega\Lambda\Omega^{-1},\quad \Omega=[x_1,\dots,x_m],\ \Lambda=\mathrm{diag}(\omega_i),
  \]
  solves it, and is stable iff $\max|\omega_i|<1$.
\end{enumerate}
\end{column}
\begin{column}{0.41\textwidth}
\begin{conceptbox}[title=\textbf{Note the symbol}]\small
The generalized eigenvalues are written $\omega$ here. The literature writes $\lambda$; this course reserves $\lambda$ for the invariant distribution, and $\chi$ is already the perturbation parameter.
\end{conceptbox}
\medskip
\begin{examplebox}\footnotesize
Quadratic in $P$, but \emph{linear} in the pencil $(\Xi,\Delta)$ --- which is exactly the trick that turns a hard nonlinear problem into a library call.
\end{examplebox}
\end{column}
\end{columns}
\vspace{1mm}
\src{FV\&G, \emph{Solution Methods} V (Appendix), ``How to solve quadratic equations''. In code: \texttt{scipy.linalg.ordqz} --- the generalized Schur (QZ) decomposition, which is Klein's (2000) route.}
\end{frame}

%=============================================================================
\begin{frame}{Blanchard--Kahn: The Condition}
\begin{columns}[T]
\begin{column}{0.55\textwidth}
\small
Write the linearized model as
\[
  A\,x_{t+1|t} = B\,x_t + C z_t,
  \qquad z_{t+1}=\rho z_t+\varepsilon_{t+1},
\]
where $x_t$ stacks states and controls. Split the variables:
\medskip
\begin{itemize}\small
  \item \textbf{Predetermined} --- capital. Cannot jump; its value today was decided yesterday.
  \item \textbf{Non-predetermined} (``jump'') --- consumption. Chooses itself now, and can move discontinuously.
\end{itemize}
\medskip
\textbf{Blanchard--Kahn condition.} There is a unique stable solution if and only if
\[
  \#\{\,|\omega_i|>1\,\} \;=\; \#\{\text{jump variables}\} .
\]
\end{column}
\begin{column}{0.42\textwidth}
\begin{conceptbox}[title=\textbf{The intuition}]\small
Each unstable root is a direction along which the economy would explode. Each jump variable is one instrument for stopping it --- choose today's value to sit on the stable manifold.
\medskip
Too few instruments and paths explode; too many and the extra ones are free.
\end{conceptbox}
\end{column}
\end{columns}
\vspace{1mm}
\src{Blanchard \& Kahn (1980), \emph{Econometrica} 48(5). Frame after Quant\_Macro Lec.~9, ``The Blanchard--Kahn conditions''.}
\end{frame}

%=============================================================================
\begin{frame}{The Three Cases}
\small
\begin{center}
\begin{tabularx}{0.97\linewidth}{@{}>{\raggedright\arraybackslash}p{2.9cm} >{\raggedright\arraybackslash}p{3.0cm} >{\raggedright\arraybackslash}X@{}}
\toprule
\textbf{Count} & \textbf{Name} & \textbf{What it means} \\
\midrule
unstable roots $=$ jumps & \textbf{Uniqueness} & Exactly one bounded path from any initial state. The normal case, and what you hope for. \\
unstable roots $<$ jumps & \textbf{Indeterminacy} & A continuum of bounded paths. Expectations can be self-fulfilling; sunspots matter. \\
unstable roots $>$ jumps & \textbf{Non-existence} & No bounded path at all. The model, as written, has no equilibrium of this kind. \\
\bottomrule
\end{tabularx}
\end{center}
\medskip
\begin{columns}[T]
\begin{column}{0.49\textwidth}
\begin{conceptbox}[title=\textbf{In the growth model}]\footnotesize
One predetermined variable ($k$), one jump ($c$). So we need \textbf{exactly one} root outside the unit circle --- which is the saddle-path property, restated as a count.
\end{conceptbox}
\end{column}
\begin{column}{0.47\textwidth}
\begin{cautionbox}\footnotesize
Indeterminacy is not a bug to be tuned away. 2.B2 met it in an OLG economy (Feng \& Hoelle 2017), where it is the substantive result.
\end{cautionbox}
\end{column}
\end{columns}
\vspace{1mm}
\src{Quant\_Macro Lec.~9, ``Interpreting Blanchard--Kahn Conditions''.}
\end{frame}

%=============================================================================
\begin{frame}{Measured: Our Two Growth Models Pass}
\begin{columns}[T]
\begin{column}{0.53\textwidth}
\small
Generalized eigenvalue moduli, computed by QZ:
\medskip
\begin{center}\footnotesize
\begin{tabularx}{\linewidth}{@{}>{\raggedright\arraybackslash}p{2.0cm} >{\raggedright\arraybackslash}X r@{}}
\toprule
\textbf{Model} & \textbf{$|\omega|$} & \textbf{Verdict} \\
\midrule
quarterly, $\delta=0.025$ & $0.950,\;0.962,\;\mathbf{1.050}$ & unique \\
Brock--Mirman, $\delta=1$ & $0.360,\;0.950,\;\mathbf{2.924}$ & unique \\
\bottomrule
\end{tabularx}
\end{center}
\medskip
One root outside the unit circle in each, one jump variable in each. The condition holds.
\medskip

The unstable root of the $\delta=1$ model is $2.923977$, and $1/(\alpha\beta)=2.923977$ --- the closed form again, arriving through an eigenvalue solver that knows nothing about it.
\end{column}
\begin{column}{0.44\textwidth}
\begin{takeawaybox}\footnotesize
Note the two stable roots: $\rho=0.95$, which is just the shock, and $0.962$, which is the endogenous persistence of capital. The model's memory is \emph{longer} than the shock's.
\end{takeawaybox}
\medskip
\begin{examplebox}\footnotesize
The unstable root here is $1.0499$, \textbf{not} $1/\beta=1.0101$. The two coincide only at $\delta=1$.
\end{examplebox}
\end{column}
\end{columns}
\vspace{1mm}
\src{Ledger \texttt{blanchard\_kahn.growth\_models} and \texttt{pert\_first\_order}.}
\end{frame}

%=============================================================================
\begin{frame}{Measured: A Model Where the Count Actually Bites}
\begin{columns}[T]
\begin{column}{0.5\textwidth}
\small
The growth model can only ever show case 1. To see the other two, take the three-equation New Keynesian model,
\begin{align*}
 \pi_t &= \beta\mathbb{E}_t\pi_{t+1}+\kappa x_t\\
 x_t &= \mathbb{E}_t x_{t+1}-\tfrac{1}{\varsigma}\big(i_t-\mathbb{E}_t\pi_{t+1}\big)\\
 i_t &= \phi_\pi \pi_t
\end{align*}
Two jump variables $(\pi,x)$, one exogenous state. So uniqueness needs \textbf{two} unstable roots.
\medskip

Sweeping the policy rule:
\end{column}
\begin{column}{0.47\textwidth}
\begin{center}\footnotesize
\begin{tabularx}{\linewidth}{@{}r >{\raggedright\arraybackslash}X r@{}}
\toprule
$\phi_\pi$ & $|\omega|$ & unstable \\
\midrule
$0.0$ & $0.704,\,0.800,\,1.435$ & 1 \\
$0.8$ & $0.800,\,0.895,\,1.244$ & 1 \\
$1.0$ & $0.800,\,1.000,\,1.139$ & 1 \\
$1.5$ & $0.800,\,1.097,\,1.097$ & \textbf{2} \\
$2.5$ & $0.800,\,1.154,\,1.154$ & \textbf{2} \\
\bottomrule
\end{tabularx}
\end{center}
\medskip
\begin{takeawaybox}\footnotesize
The switch happens at $\phi_\pi=1$: the \textbf{Taylor principle}, recovered from an eigenvalue count. Below it the equilibrium is indeterminate; at exactly $1$ a root sits on the unit circle.
\end{takeawaybox}
\end{column}
\end{columns}
\vspace{1mm}
\src{Ledger \texttt{blanchard\_kahn.new\_keynesian}: $\kappa=0.1275$, $\varsigma=1$, $\rho_u=0.8$. For $\phi_\pi<1$ the QZ routine cannot even form $P$ --- the block it must invert is singular, which is what indeterminacy looks like from inside the algorithm.}
\end{frame}

%=============================================================================
\begin{frame}{In Practice: Who Does This For You}
\begin{columns}[T]
\begin{column}{0.53\textwidth}
\small
Four procedures solve the quadratic system, and they are equivalent: \textbf{Blanchard \& Kahn (1980)}, \textbf{Uhlig (1999)}, \textbf{Sims (2002)}, \textbf{Klein (2000)}. They differ in what they assume about the matrices, not in the answer.
\medskip

In practice you use software: \textbf{Dynare} for first, second and third order; Dynare++ or Mathematica beyond that. You supply the equilibrium conditions; it does the steady state, the derivatives and the root selection.
\medskip

The real burden is the \textbf{derivatives}. Dynare computes them analytically by symbolic differentiation, and for a reason.
\end{column}
\begin{column}{0.44\textwidth}
\begin{cautionbox}[title=\textbf{Why not numerical derivatives?}]\small
We used them, and watched them fail. First order: the finite-difference solve matched QZ to $4.4\times10^{-11}$. Second order, with a step ten times smaller, the residual of the same system went from $3.6\times10^{-9}$ to $\mathbf{8\times10^{-3}}$.
\medskip
Second derivatives divide by $h^{2}$, so rounding error is amplified by $10^{8}$ --- 4.A4's U-curve, in a place where it silently corrupts a published coefficient.
\end{cautionbox}
\end{column}
\end{columns}
\vspace{1mm}
\src{FV\&G, \emph{Solution Methods} V, ``Solving the system II'' and ``Computer''. The failure above is from this deck's own build log, not a source. Automatic differentiation (4.A4) is the modern answer.}
\end{frame}

%=============================================================================
\begin{frame}{Takeaways}
\begin{takeawaybox}
\begin{itemize}\small
  \item \textbf{Log-linearize} by default: coefficients are elasticities, deviations are percentages --- and for this model it is exact.
  \item Guess linear rules, substitute, match coefficients: the state block collapses to a \textbf{matrix quadratic} in $P$, stable iff $\max|\mathrm{eig}(P)|<1$. It is quadratic because the model has a \textbf{stable and an unstable manifold} --- discarding the explosive root \emph{is} imposing transversality.
  \item Solve it as a \textbf{generalized eigenvalue problem}: the QZ decomposition, one library call, Klein's (2000) route.
  \item \textbf{Blanchard--Kahn:} unique iff (\#roots outside the unit circle) $=$ (\#jump variables). Fewer gives indeterminacy, more gives non-existence.
  \item Measured: both growth models have exactly one unstable root ($1.050$ quarterly, $2.924=1/(\alpha\beta)$ at $\delta=1$); in the New Keynesian model the count flips at $\phi_\pi=1$ --- \textbf{the Taylor principle as an eigenvalue count}.
  \item Use Dynare, and let it differentiate \textbf{analytically}: our own numerical second derivatives cost six orders of magnitude of residual.
\end{itemize}
\end{takeawaybox}
\end{frame}

%=============================================================================
\begin{frame}{Bridge: What the Linear Solution Cannot See}
\begin{columns}[T]
\begin{column}{0.53\textwidth}
\small
We now have a complete, unique, checkable first-order solution. It also has, by construction, \textbf{no risk in it at all}: 5.A3 showed $c_\chi=k_\chi=0$, so the linear rules are the same whether the shock has standard deviation $0.007$ or $0$.
\medskip

For dynamics that is often fine. For anything about uncertainty --- precautionary saving, risk premia, the welfare cost of business cycles --- it returns exactly zero, which is not an estimate.
\end{column}
\begin{column}{0.44\textwidth}
\begin{conceptbox}[title=\textbf{Next (5.A5)}]\small
Second order: the twelve-equation system, the risk correction, and why every odd derivative in $\chi$ vanishes. Then why second-order simulations explode, and the pruning that stops them.
\end{conceptbox}
\medskip
\begin{examplebox}\footnotesize
Also: how far from the steady state any of this is valid, which turns out to be a question about \emph{complex} analysis.
\end{examplebox}
\end{column}
\end{columns}
\end{frame}

\end{document}
